% Complete independent mathematical fragment HPS. Requires amsmath,amssymb,hyperref.
\section{Exact prime dilation of the original shifted Hurwitz family}
\label{sec:HPS}

\paragraph{Source, scope, and coordinates.}
The supplied programme source is
\texttt{split\_zero\_cue\_shifted\_hurwitz\_phase\_space.tex},
SHA--256
\begin{center}\small
\texttt{52f2a6bccdb15213b1992abef9f8d6476d6e9ea5566ee7ed9d85c9a2d7cc2705}.
\end{center}
The actual source reading for this calculation is lines
250--350 and 545--686, including the complete proofs of the
split-zero laws, support reflection, universal shifted-flow
equation, and non-Eulerian shifted sheet. The source locators
are \texttt{thm:universal-flow} and \texttt{prop:non-euler}.
The full-lattice laws and synchronized lifts are restated
from FLH1--3; their complete proofs were written and read
in the preceding programme calculation. The tropical
coefficient embedding is SHP15--16, directly read here
in lines 330--401 of \texttt{SUPPORTED\_HEAT\_PRODUCT.tex},
at SHA--256
\begin{center}\small
\texttt{0b0f5a50874ce59fdfdb74a715881182a38a9238c816107d0d3b143403426799}.
\end{center}
Every operator and convergence assertion needed below is
proved here. No identification with a Bost--Connes system
or assertion about zero locations is made.

For the exact classical comparison, Alain Connes,
Caterina Consani and Matilde Marcolli,
\href{https://arxiv.org/abs/0806.2401v1}
{\emph{Fun with $\mathbb F_1$}}, original author file
\texttt{funBC.tex}, lines 804--885, was read directly,
including its equation \texttt{sigmatgenerators}.
Its stated dilation evolution is the scalar
$m^{iu}$ law at zero deformation, with the
source's time renamed $u$ only to distinguish
it from the original deformation $t$.
The file has SHA--256
\begin{center}\small
\texttt{ab76adc442f748a75eb5cb0a099fcbb4dcbd96f8cbc7055efca846b79a262b48}.
\end{center}
Eugene Ha and Fr\'ed\'eric Paugam,
\href{https://arxiv.org/abs/math/0507101v1}
{\emph{Bost--Connes--Marcolli systems for Shimura
varieties. I. Definitions and formal analytic properties}},
original author source lines 461--504, was also
read directly: its Section 1.1 gives $H=\log N$
and $\operatorname{Tr}(e^{-sH})=\zeta(s)$.
The source SHA--256 is
\begin{center}\small
\texttt{7efe5e2c38c48aefaed921db505b1454d36873913617be28bf86b92ec738bf5c}.
\end{center}
These are bounded source readings. The shifted
relations below are calculated on the original
Hurwitz family; no larger identification is
deduced from those two classical formulas.

Throughout, $t>-1$ is the original additive deformation,
$u\in\mathbb R$ is the unitary evolution time, and
$s\in\mathbb C$ is the trace exponent, with
$\sigma=\operatorname{Re}s>1$ wherever a partition trace
or an unrestricted infinite prime product is used.
The vanishing-at-zero subspace in HPS33--39 has
the explicitly proved larger domain $\operatorname{Re}s>0$.
These three parameters
are not identified.

\subsection{Unbounded operators and the original partition trace}

On $\mathcal H=\ell^2(\mathbb N_{\ge1})$ with its original
orthonormal basis $(e_n)_{n\ge1}$ define
\begin{equation}\tag{HPS1}
\begin{gathered}
 Ne_n=ne_n,\qquad
 \mathcal D(N)=\{x:\sum_{n\ge1}n^2|x_n|^2<\infty\},\\
 H_te_n=\log(n+t)e_n,\qquad
 \mathcal D(H_t)=\{x:\sum_{n\ge1}|\log(n+t)|^2|x_n|^2<\infty\}.
\end{gathered}
\end{equation}
All logarithms of positive real numbers are real.
These multiplication operators are self-adjoint:
finite sequences lie in their domains and are dense.
If $y$ belongs to the adjoint domain of a real diagonal
operator with entries $a_n$, testing against $e_n$
forces the adjoint coordinates to be $a_ny_n$; they
must be square summable. Conversely that condition
makes the inner-product identity valid by
Cauchy--Schwarz. Thus the adjoint has exactly the
stated domain and values. In particular $N+tI$ is
strictly positive, with lower bound $1+t$, and its
spectral logarithm is exactly $H_t$.

There is one common logarithmic domain:
\begin{equation}\tag{HPS2}
\begin{gathered}
 \mathcal D=\{x:\sum_{n\ge1}(\log n)^2|x_n|^2<\infty\},
 \qquad \mathcal D(H_t)=\mathcal D,\\
 H_t=H_0+B_t,\qquad
 B_te_n=\log(1+t/n)e_n,\qquad
 \|B_t\|=|\log(1+t)| .
\end{gathered}
\end{equation}
Indeed the diagonal difference is bounded, with its
largest absolute value at $n=1$, and tends to zero.
The two inequalities
$|a+b|^2\le2|a|^2+2|b|^2$ applied in both directions
show equality of the domains. This also proves that
$B_t$ is compact: truncating its diagonal after $M$
has error the supremum of its remaining entries.

Define complex powers by their actual diagonal entries.
For $\sigma>1$,
\begin{equation}\tag{HPS3}
\begin{gathered}
 A_t(s)=e^{-sH_t},\qquad A_t(s)e_n=(n+t)^{-s}e_n,\\
 \|A_t(s)\|_1=\sum_{n\ge1}(n+t)^{-\sigma}<\infty,\qquad
 \operatorname{Tr}A_t(s)
       =F_t(s)=\sum_{n\ge1}(n+t)^{-s}=\zeta(s,1+t).
\end{gathered}
\end{equation}
For a diagonal operator its singular values are the
absolute diagonal entries, with multiplicities: the
positive square root of $A^*A$ is the diagonal of
their absolute values. The trace-class criterion and
trace formula therefore follow by finite truncation.
Comparison with $n^{-\sigma}$ proves convergence
exactly for $\sigma>1$; at $\sigma\le1$ the absolute
diagonal sum diverges. For $\sigma<0$ the diagonal
power is unbounded, so in that case it is not a bounded
trace-class operator in the first place.
For real $\beta>1$, $A_t(\beta)$ is positive, and
$A_t(\beta)/F_t(\beta)$ is a positive trace-one
operator. The denominator is retained in this
explicit definition; $F_t$ itself is unchanged.

The original differential flow is an actual convergent
trace identity:
\begin{equation}\tag{HPS4}
\begin{gathered}
 \partial_t^jF_t(s)=(-1)^j(s)_jF_t(s+j),\qquad
 (s)_j=s(s+1)\cdots(s+j-1),\quad(s)_0=1,\\
 \partial_tF_t=-DF_t,\qquad (Df)(s)=sf(s+1),\\
 F_{t+r}(s)=\sum_{j\ge0}
       \frac{(-r)^j(s)_j}{j!}F_t(s+j)
                         \qquad(|r|<1+t).
\end{gathered}
\end{equation}
For a real compact interval of $t>-1$,
$n+t\ge c n$ for some $c>0$. On a compact set of
$s$ with real part greater than one, termwise
derivatives have a summable bound
$C_j n^{-\sigma_0-j}$, and derivatives in $s$ add
only fixed powers of $\log(n+t)$.
One may also let $t$ be complex in $\operatorname{Re}t>-1$:
every $n+t$ lies in the right half-plane and uses
its principal logarithm; the same compact bounds
give holomorphy by uniform power-series convergence.
For the last identity fix $|r|<R<1+t$ and apply the
binomial Taylor series to each $(n+t+r)^{-s}$.
Its absolute sum is bounded uniformly in $n$ by
$C_{s,R,r}(n+t)^{-\sigma}$: the Taylor coefficients
of $(1+w)^{-s}$ are bounded on $|w|=R/(1+t)<1$,
and $|r|/(n+t)\le|r|/(1+t)$.
The geometric Cauchy bound sums absolutely.
This proves every exchange and the exact radius
asserted in HPS4. It does not claim a globally
convergent Taylor series across $t=-1$.

\subsection{Exact dilation, its fixed-parameter defect, and coherence}

For each positive integer $m$ let $S_me_n=e_{mn}$.
It is an isometry, with
$S_m^*e_j=e_{j/m}$ if $m\mid j$ and $0$ otherwise,
and $S_mS_k=S_{mk}=S_kS_m$.
The original unbounded operators satisfy
\begin{equation}\tag{HPS5}
\begin{gathered}
 x\in\mathcal D\ \Longleftrightarrow\ S_mx\in\mathcal D,\qquad
 S_m^*\mathcal D\subseteq\mathcal D,\\
 H_tS_mx=S_m(H_{t/m}+\log m\,I)x
                         \quad(x\in\mathcal D),\\
 S_m^*H_tS_m=H_{t/m}+\log m\,I
                         \quad\hbox{on domain }\mathcal D.
\end{gathered}
\end{equation}
The scalar identity is
$\log(mn+t)=\log m+\log(n+t/m)$.
Also $\log(mn)=\log m+\log n$, so square-summability
with the logarithmic weight is equivalent before
and after $S_m$; for $S_m^*$ use
$\log n\le\log(mn)$ on the selected coordinates.
Thus these are equalities with their full domains,
not merely identities on basis vectors.

Let $U_t(u)e_n=e^{iu\log(n+t)}e_n$.
The diagonal modulus is one, so $U_t$ is a unitary
group; dominated convergence in
$\sum |x_n|^2$ proves strong continuity.
Exponentiating HPS5 coordinatewise gives the exact
parameter-changing scaling
\begin{equation}\tag{HPS6}
\begin{gathered}
 U_t(u)S_mU_{t/m}(-u)=m^{iu}S_m,\\
 A_t(s)S_m=m^{-s}S_mA_{t/m}(s),\qquad
 S_m^*A_t(s)S_m=m^{-s}A_{t/m}(s)
                                  \quad(\sigma>1).
\end{gathered}
\end{equation}
Thus $S_m$ from the copy carrying $H_{t/m}$ to the
copy carrying $H_t$ restores exact scaling.
For the composite dilation, $S_k$ first carries
$H_{t/(mk)}$ to $H_{t/m}$, and $S_m$ carries that
copy to $H_t$. The additive energy shifts are
$\log k+\log m=\log(mk)$. Hence this changing of
parameters is coherent for every finite product.

At a fixed deformation define
$\alpha_u^t(B)=U_t(u)BU_t(-u)$ on
$\mathcal B(\mathcal H)$. Direct calculation gives
\begin{equation}\tag{HPS7}
\begin{gathered}
 \alpha_u^t(S_m)=S_mC_m(t,u;N),\\
 C_m(t,u;n)=
       \exp\!\left(iu\log\frac{mn+t}{n+t}\right)
       =m^{iu}\exp(iu b_m(t;n)),\\
 b_m(t;n)=\log(n+t/m)-\log(n+t).
\end{gathered}
\end{equation}
These are bounded unitary diagonal operators, not
constant coefficients in general. If $m>1$, $t\ne0$,
and $u\ne0$, equality $C_m(t,u;N)=m^{iu}I$
would require $u b_m(t;n)\in2\pi\mathbb Z$ for all
$n$. For sufficiently large $n$ this number is
nonzero and has modulus less than $2\pi$, a
contradiction. The fixed-parameter scalar formula
therefore fails exactly in this nontrivial case.

Both the time cocycle and the multiplicative
coherence are explicit:
\begin{equation}\tag{HPS8}
\begin{gathered}
 C_m(t,u+v;n)=C_m(t,u;n)C_m(t,v;n),\\
 C_m(t,u;kn)C_k(t,u;n)=C_{mk}(t,u;n)
                    =C_k(t,u;mn)C_m(t,u;n),\\
 b_m(t;kn)+b_k(t;n)=b_{mk}(t;n)
                    =b_k(t;mn)+b_m(t;n).
\end{gathered}
\end{equation}
The ratios in the second line telescope.
Subtracting the real constants $\log m,\log k$
proves the third line, with no logarithm branch
ambiguity. To verify the operator multiplication,
use $f(N)S_k=S_kf(kN)$ on each coordinate.
It gives
$\alpha_u^t(S_m)\alpha_u^t(S_k)
=S_{mk}C_m(t,u;kN)C_k(t,u;N)$,
which is exactly $\alpha_u^t(S_{mk})$.

The defect from scalar scaling has an evaluated
operator size. Put $c_t=\min(t,0)>-1$. Then
\begin{equation}\tag{HPS9}
\begin{gathered}
 |b_m(t;n)|\le
       \frac{|t|(1-1/m)}{n+c_t},\qquad
 \lim_{n\to\infty}n b_m(t;n)=-t(1-1/m),\\
 R_{m,t,u}:=\alpha_u^t(S_m)-m^{iu}S_m,\\
 \|R_{m,t,u}\|_{\mathcal S^q}^{\,q}
 \le |ut|^q(1-1/m)^q\sum_{n\ge1}(n+c_t)^{-q}
                  <\infty\qquad(q>1),\\
 t u\ne0,\ m>1\quad\Longrightarrow\quad
       R_{m,t,u}\notin\mathcal S^1 .
\end{gathered}
\end{equation}
Indeed $b_m$ is the oriented integral of
$(n+v)^{-1}$ between $t$ and $t/m$.
The denominator on that segment is at least $n+c_t$,
which proves the bound; multiplication by $n$ and
uniform convergence of $n/(n+v)$ on the fixed
segment proves the limit. Since $S_m$ is an
isometry, $R^*R$ is diagonal with entries
$|e^{iu b_m(t;n)}-1|^2$. These are its squared
singular values, whether or not written in
decreasing order. The inequality
$|e^{iy}-1|\le|y|$ proves the Schatten bound and
compactness. In the nontrivial case their square
roots satisfy
$n|e^{iu b_m(t;n)}-1|\to|ut|(1-1/m)>0$.
The harmonic series proves failure of trace class.
Thus an actual compact correction has been
calculated; it has not been discarded.

\subsection{Divisibility projections and the exact prime strata}

Let $P_m=S_mS_m^*$, the orthogonal projection onto
the multiples of $m$. All such projections are
diagonal and
\begin{equation}\tag{HPS10}
 P_mP_k=P_{\operatorname{lcm}(m,k)},\qquad
 P_mA_t(s)=A_t(s)P_m,\qquad
 \operatorname{Tr}(P_mA_t(s))=m^{-s}F_{t/m}(s).
\end{equation}
The first identity is the intersection of two
divisibility conditions. The trace follows by
writing its selected index as $mn$ and using
$mn+t=m(n+t/m)$. It is also the compression
identity HPS6; absolute trace convergence makes
both computations agree.

For a finite set $\mathcal P$ of primes let
$M=\prod_{p\in\mathcal P}p$, and let $\mu(d)$ be
$(-1)^r$ for a product of $r$ distinct primes,
$0$ for an integer divisible by a prime square,
and $1$ at $d=1$. Put
\begin{equation}\tag{HPS11}
\begin{gathered}
 Q_{\mathcal P}=\prod_{p\in\mathcal P}(I-P_p),\\
 R_{\mathcal P}(t;s)=\operatorname{Tr}
                   (Q_{\mathcal P}A_t(s))
   =\sum_{\substack{n\ge1\\(n,M)=1}}(n+t)^{-s}
   =\sum_{d\mid M}\mu(d)d^{-s}F_{t/d}(s).
\end{gathered}
\end{equation}
All products are finite. Their expansion gives
$\sum_{d\mid M}\mu(d)P_d$, since distinct-prime
least common multiples are products. On each
integer the result is one precisely when no
prime in $\mathcal P$ divides it, and zero
otherwise. HPS10 proves the final trace
identity with every deformation $t/d$ retained.

For prescribed valuations $a_p\ge0$ at
$p\in\mathcal P$, write $d=\prod p^{a_p}$.
The exact stratum projection and trace are
\begin{equation}\tag{HPS12}
\begin{gathered}
 E_{\mathcal P,\boldsymbol a}
   =\prod_{p\in\mathcal P}
                (P_{p^{a_p}}-P_{p^{a_p+1}})
   =S_dQ_{\mathcal P}S_d^*,\\
 \operatorname{Tr}(E_{\mathcal P,\boldsymbol a}A_t(s))
        =d^{-s}R_{\mathcal P}(t/d;s)
        =\sum_{e\mid M}\mu(e)(de)^{-s}F_{t/(de)}(s),\\
 \sum_{\boldsymbol a\in\mathbb N_0^{\mathcal P}}
       E_{\mathcal P,\boldsymbol a}=I
                     \quad\hbox{strongly},\qquad
 F_t(s)=\sum_{\boldsymbol a}
                    d^{-s}R_{\mathcal P}(t/d;s).
\end{gathered}
\end{equation}
Every coordinate integer has exactly one valuation
vector; the projections therefore are orthogonal,
and their finite partial sums converge strongly by
the square-summable coordinate tail.
Writing the stratum integer as $dr$ with $(r,M)=1$
proves its trace. The absolute sum over all strata
is at most $\sum_n(n+t)^{-\sigma}$, so the last
trace sum has no conditional rearrangement.

There is also an additive decomposition using
every prime exactly as a least prime divisor.
Let $\mathcal P_{<p}$ contain all primes less
than $p$, and let $\Pi_1$ project onto $\mathbb C e_1$.
Then
\begin{equation}\tag{HPS13}
\begin{gathered}
 I=\Pi_1+\sum_{p\ {\rm prime}}P_pQ_{\mathcal P_{<p}}
                                \quad\hbox{strongly},\\
 F_t(s)=(1+t)^{-s}
       +\sum_{p\ {\rm prime}}p^{-s}
                    R_{\mathcal P_{<p}}(t/p;s).
\end{gathered}
\end{equation}
For $n>1$, its least prime divisor is unique;
the indicated projections select exactly these
disjoint sets. For $n=pr$, having no smaller
prime divisor is equivalent to $(r,\prod_{q<p}q)=1$.
This proves the second line by the same absolutely
convergent trace partition as HPS12.

\subsection{A convergent operator Euler product on the full parameter family}

Fix a compact interval $J=[a,b]\subset(-1,\infty)$
containing $0$; every compact set of allowed real
deformations is contained in such a $J$.
Let $X_J=C(J,\mathbb C)$ with the supremum norm
and pointwise ring operations. For $m\ge1$ put
\begin{equation}\tag{HPS14}
 (V_mf)(t)=f(t/m),\qquad
 K_m(s)=m^{-s}V_m,\qquad
 V_mV_n=V_{mn},\quad K_m(s)K_n(s)=K_{mn}(s),\quad
 \|K_m(s)\|=m^{-\sigma}.
\end{equation}
The interval is stable under division by $m$.
The norm upper bound follows from the supremum;
the constant function $1$ attains it. The
composition formulas are literal substitutions.
The maps $V_m$ are unital ring homomorphisms;
$K_m$ are bounded linear operators, and generally
are not ring homomorphisms because of their
scalar factor.

Define bounded linear operators by absolutely
convergent operator series:
\begin{equation}\tag{HPS15}
\begin{gathered}
 \mathcal E_s=\sum_{n\ge1}n^{-s}V_n,\qquad
 \mathcal M_s=\sum_{n\ge1}\mu(n)n^{-s}V_n,\\
 \|\mathcal E_s\|\le\zeta(\sigma),\qquad
 \|\mathcal M_s\|\le\sum_n|\mu(n)|n^{-\sigma}
                                      \le\zeta(\sigma),\\
 \mathcal M_s\mathcal E_s=\mathcal E_s\mathcal M_s=I.
\end{gathered}
\end{equation}
Both absolute norm sums are bounded by
$\sum n^{-\sigma}$. Their double product may
therefore be regrouped by the product index:
the coefficient of $n^{-s}V_n$ is
$\sum_{d\mid n}\mu(d)$. If $n$ has $r$ distinct
prime factors this is $(1-1)^r$, hence it is
$0$ for $n>1$ and $1$ for $n=1$. This proves
both inverse identities, using only finite
prime factorization and the proved absolute
operator sums.

For a finite prime set $\mathcal P$ write
$\mathbb N_{\mathcal P}$ for its smooth positive
integers, that is, those whose prime factors
belong to $\mathcal P$. The exact finite products are
\begin{equation}\tag{HPS16}
\begin{gathered}
 \mathcal E_{s,\mathcal P}
   =\prod_{p\in\mathcal P}(I-K_p(s))^{-1}
   =\prod_{p\in\mathcal P}\sum_{j\ge0}K_{p^j}(s)
   =\sum_{n\in\mathbb N_{\mathcal P}}n^{-s}V_n,\\
 \mathcal M_{s,\mathcal P}
   =\prod_{p\in\mathcal P}(I-K_p(s))
   =\sum_{d\mid M}\mu(d)d^{-s}V_d,\\
 \lim_{\mathcal P}\mathcal E_{s,\mathcal P}
       =\mathcal E_s,\qquad
 \lim_{\mathcal P}\mathcal M_{s,\mathcal P}
       =\mathcal M_s
                           \quad\hbox{in operator norm}.
\end{gathered}
\end{equation}
Since $\|K_p\|=p^{-\sigma}<1$, the geometric
series converges in operator norm and its
finite telescoping product with $I-K_p$
has remainder $K_p^{j+1}\to0$. It is therefore
the actual inverse. Commutativity and unique
prime factorization prove the first expansion;
its multiple geometric sum is absolute.
The second product is finite inclusion-exclusion.
As the finite prime set increases, each integer
eventually occurs in the first expansion, and
each square-free integer in the second.
Dominating their missing coefficient norms by
the summable $n^{-\sigma}$ proves both limits.
These are inverses in a Banach algebra of linear
operators; no subtraction in a split-zero
semiring has been used to define an inverse.

The convergence has the following uniform bound.
If $\mathcal P$ contains every prime at most
$x\ge2$ and $q=\lfloor x\rfloor$, then
\begin{equation}\tag{HPS17}
 \|\mathcal E_s-\mathcal E_{s,\mathcal P}\|,
 \ \|\mathcal M_s-\mathcal M_{s,\mathcal P}\|
 \ \le\sum_{n>q}n^{-\sigma}
 \ \le\frac{q^{1-\sigma}}{\sigma-1}.
\end{equation}
An omitted integer has a prime divisor greater
than $x$, hence is greater than $q$.
The decreasing function $v^{-\sigma}$ gives
the integral bound from $q$ to infinity.
On compact sets with $\operatorname{Re}s\ge
1+\epsilon$, the same estimate and all
operator sums are uniform; differentiating
in $s$ adds $(\log n)^j$, still summable
on each such compact set. Thus these are
holomorphic operator families there.

The original shifted trace is obtained by applying
this actual Euler product to its first atom:
\begin{equation}\tag{HPS18}
\begin{gathered}
 g_s(t)=(1+t)^{-s},\qquad Z_s(t)=F_t(s),\\
 Z_s=\mathcal E_sg_s
     =\lim_{\mathcal P}
             \prod_{p\in\mathcal P}(I-K_p(s))^{-1}g_s,
 \qquad
 \mathcal M_sZ_s=g_s,\\
 (\mathcal M_{s,\mathcal P}Z_s)(t)
           =R_{\mathcal P}(t;s).
\end{gathered}
\end{equation}
Indeed each term is exactly
$n^{-s}(1+t/n)^{-s}=(n+t)^{-s}$, because the
bases are positive and real for $t\in J$.
The series is uniform on $J$ by HPS15,
and it agrees term for term with HPS3.
The inverse formula is HPS15; the finite
sieve is HPS11. The errors in the two
finite-product formulas are bounded by
the right side of HPS17 times respectively
$\|g_s\|_\infty$ and $\|Z_s\|_\infty$.
This proves local uniform convergence in the
original parameter, not only pointwise convergence.

Fixed-parameter evaluation retains the exact
reason for the changing prime factors:
\begin{equation}\tag{HPS19}
\begin{gathered}
 \operatorname{ev}_tK_m(s)
                =m^{-s}\operatorname{ev}_{t/m},\\
 \operatorname{ev}_0\mathcal E_s
                =\zeta(s)\operatorname{ev}_0,\qquad
 F_0(s)=\zeta(s),\qquad g_s(0)=1.
\end{gathered}
\end{equation}
For $t\ne0$ and $m>1$, evaluation at $t/m$
and at $t$ are different functionals on $X_J$:
the test function $f(v)=v$ distinguishes them.
For the actual atom, their modulus differs
as well, since $\sigma>1$ and
$1+t/m\ne1+t$.
Thus no scalar Euler factor at the original
fixed $t$ replaces this parameter-changing
operator. At $t=0$, every $V_m$ acts as the
identity after evaluation; HPS16 returns
the usual scalar Euler product directly.

The operator prime product also respects the
original shifted differential equation.
For jointly differentiable families $f(s,t)$ set
$(Df)(s,t)=sf(s+1,t)$. Then
\begin{equation}\tag{HPS20}
\begin{gathered}
 \partial_t(K_m(s)f)
           =m^{-1}K_m(s)\partial_tf,\qquad
 D(K_m(s)f)=m^{-1}K_m(s)Df,\\
 (\partial_t+D)(K_m(s)f)
           =m^{-1}K_m(s)(\partial_t+D)f,\\
 \partial_t(\mathcal E_sf)
            =\mathcal E_{s+1}\partial_tf
                         \qquad(f\in C^1(J,\mathbb C)).
\end{gathered}
\end{equation}
The first formula is the chain rule. In the
second, evaluating the scalar factor at $s+1$
gives $s m^{-(s+1)}f(s+1,t/m)$, exactly the
displayed right side. The atom $g_s(t)$
satisfies $(\partial_t+D)g=0$.
For the last formula take a nondegenerate
interval $J$ and use its usual one-sided
endpoint derivatives. The differentiated
operator series is uniformly bounded by
$\sum n^{-\sigma-1}\|f'\|_\infty$.
The original series is uniformly bounded by
$\sum n^{-\sigma}\|f\|_\infty$.
The fundamental theorem of calculus on $J$
therefore proves that its sum is $C^1$
and that its derivative is the displayed
sum. In particular
$\|\mathcal E_sf\|_{C^1}
\le\zeta(\sigma)\|f\|_\infty
+\zeta(\sigma+1)\|f'\|_\infty$
for the sum norm. Applying the formula to
$g_s$ gives
$\partial_tZ_s=-s\mathcal E_{s+1}g_{s+1}
=-s Z_{s+1}$ with the exact original
operator $D$. Thus HPS20 recovers the
original Hurwitz flow from the convergent
prime-factor construction.

For completeness the original Gibbs divisibility
correlation is also evaluated. For real $\beta>1$
and distinct primes $p,q$, write expectations with
the explicitly defined state from HPS3. Then
\begin{equation}\tag{HPS21}
\begin{gathered}
 \mathbb E_{t,\beta}P_p
     =p^{-\beta}\frac{F_{t/p}(\beta)}{F_t(\beta)},\\
 \operatorname{Cov}_{t,\beta}(P_p,P_q)
  =(pq)^{-\beta}
       \frac{F_t(\beta)F_{t/(pq)}(\beta)
                    -F_{t/p}(\beta)F_{t/q}(\beta)}
            {F_t(\beta)^2},\\
 \left.\partial_t\operatorname{Cov}_{t,\beta}(P_p,P_q)
        \right|_{t=0}
  =-(pq)^{-\beta}\beta
       \frac{\zeta(\beta+1)}{\zeta(\beta)}
                         (1-1/p)(1-1/q)<0.
\end{gathered}
\end{equation}
Use $P_pP_q=P_{pq}$ and HPS10 for the first
two lines. The numerator at $0$ vanishes.
HPS4 gives
$\partial_tF_{t/d}(\beta)|_0
=-\beta\zeta(\beta+1)/d$.
Its four product derivatives yield
$-\beta\zeta(\beta)\zeta(\beta+1)
(1+1/(pq)-1/p-1/q)$, and division by
$F_0(\beta)^2$ proves the final line.
This is a calculated failure of the undeformed
prime independence to first order, with its
original parameter and trace weights preserved.

\subsection{Every split-zero label, the coefficient tensor map, and trace}

Let $L$ be the full bounded distributive lattice,
with bottom $0_L$ and top $1_L$.
For any unital commutative complex algebra $R$,
retain the original semiring
\begin{equation}\tag{HPS22}
\begin{gathered}
 G_L(R)=\{(0,\lambda):\lambda\in L\}
                  \cup\{(r,1_L):r\in R\},\\
 (r,\lambda)+(r',\lambda')
             =(r+r',\lambda\vee\lambda'),\qquad
 (r,\lambda)(r',\lambda')
             =(rr',\lambda\wedge\lambda'),\\
 \tau=(0,0_L),\qquad z_\lambda=(0,\lambda),
 \qquad e=z_{1_L},\qquad \widehat r=(r,1_L).
\end{gathered}
\end{equation}
The same component rules define a synchronized
$G_L(\mathbb C)$-semimodule $V_L^\uparrow$ for
every complex vector space $V$.
Closure holds because every nonzero amplitude
forces top support. Associativity and all
distributive laws follow componentwise from
the ring laws and distributivity of the
lattice. The actual additive zero is $\tau$;
$\widehat0=e$, and all $z_\lambda$ are retained.

Every complex linear map $T:V\to W$ has the lift
\begin{equation}\tag{HPS23}
\begin{gathered}
 T^\uparrow(v,\lambda)=(Tv,\lambda),\\
 (TU)^\uparrow=T^\uparrow U^\uparrow,\qquad
 (T+U)^\uparrow=T^\uparrow+U^\uparrow,\qquad
 0^\uparrow(v,\lambda)=(0,\lambda).
\end{gathered}
\end{equation}
The map is defined since a nonzero output
requires a nonzero input and hence top
support. Component addition and scalar meet
prove linearity. Composition preserves the
label once, and the pointwise sum joins two
copies of the same label, proving the
identities. These statements apply also to
$H_t:\mathcal D\to\mathcal H$ with the
exact lifted domain $\mathcal D_L^\uparrow$.
Consequently HPS5--8 lift with every label
and their specified domains. The scalar
$m^{iu}$ is lifted as $\widehat{m^{iu}}$.

On $(X_J)_L^\uparrow$ the inverse maps in
HPS15--16 become exact linear isomorphisms:
\begin{equation}\tag{HPS24}
\begin{gathered}
 \mathcal M_s^\uparrow\mathcal E_s^\uparrow=I,\qquad
 \mathcal E_s^\uparrow(g_s,1_L)=(Z_s,1_L),\\
 (I-K_p)^\uparrow(f,\lambda)
       =(f-K_pf,\lambda)
       =I^\uparrow(f,\lambda)
                         +\widehat{-1}K_p^\uparrow(f,\lambda),\\
 \left(\sum_{j\ge0}K_{p^j}\right)^\uparrow(f,\lambda)
             =\left(\sum_{j\ge0}K_{p^j}f,\lambda\right).
\end{gathered}
\end{equation}
The last sum means the proved convergent
amplitude sum with its one unchanged label.
Every finite partial sum already has that
label, so no arbitrary infinite join in $L$
is required. The same interpretation applies
to all norm limits in HPS15--18.
If $f-K_pf$ vanishes, the result is
$z_\lambda$, not necessarily $\tau$.
Thus these are lifted linear inverses;
the proof does not invent additive inverses
for a split-zero semiring.

Let $K$ be the original tropical semifield
with zero and let $K_L=K\otimes_{\mathbb B}L$.
The complete coefficient map is
\begin{equation}\tag{HPS25}
\begin{gathered}
 \iota_L(\lambda)=1_K\otimes\lambda,\qquad
 \rho_L=\beta_K\otimes\operatorname{id}_L,\qquad
 \beta_K(a)=\begin{cases}0,&a=0_K,\\1,&a\ne0_K,\end{cases}\\
 \rho_L\iota_L=\operatorname{id}_L,\qquad
 \Phi_R:G_L(R)\longrightarrow R\times K_L,\quad
            (r,\lambda)\longmapsto(r,\iota_L(\lambda)).
\end{gathered}
\end{equation}
In a tropical semifield a sum is zero only
when both terms are zero, and a product is
zero exactly when a factor is zero. Hence
$\beta_K$ is a unital semiring homomorphism.
The tensor universal property gives $\rho_L$,
and its value on $1_K\otimes\lambda$ is
$\lambda$. This proves injectivity of
$\iota_L$. It preserves join, meet, zero and
top by its tensor construction, so the
component laws HPS22 prove that $\Phi_R$ is
an injective unital semiring homomorphism.
Its image is exactly the pairs with support
in $\iota_L(L)$, and with support equal to
$1_{K_L}$ whenever the amplitude is nonzero.
For synchronized vector spaces use the same
formula as an injective additive map.
For every linear $T$ the exact commuting
identity is
$(T,\operatorname{id})\Phi_V=\Phi_WT^\uparrow$.
It applies to every operator and amplitude
limit above; the retraction recovers every
original label, including $\tau$, $e$, and
all intermediate zeros.

The ordinary trace has the precisely typed lift
\begin{equation}\tag{HPS26}
\begin{gathered}
 \operatorname{Tr}^\uparrow:
       \mathcal S^1(\mathcal H)_L^\uparrow
             \longrightarrow G_L(\mathbb C),\qquad
 (A,\lambda)\longmapsto(\operatorname{Tr}A,\lambda),\\
 \Phi_{\mathbb C}\operatorname{Tr}^\uparrow
   =(\operatorname{Tr},\operatorname{id})
                    \Phi_{\mathcal S^1},\\
 \operatorname{Tr}^\uparrow(P_mA_t(s),1_L)
       =(m^{-s}F_{t/m}(s),1_L),\\
 \operatorname{Tr}^\uparrow(Q_{\mathcal P}A_t(s),1_L)
       =(R_{\mathcal P}(t;s),1_L).
\end{gathered}
\end{equation}
Trace is bounded and complex linear on its
trace-class domain, so HPS23 proves the
first two lines. The others are the
already proved amplitudes HPS10--11.
This is not a multiplicative map:
the orthogonal rank-one projections onto
$e_1,e_2$ have product zero, while both
individual traces equal one. At a vanishing
trace, the top-supported output remains $e$.

\subsection{Independent coordinate labels and the exact support defect of a sieve}

The synchronized lift can be supplemented by
every independent original coordinate label,
without identifying the two carriers. For
$r=1,2$ define
\begin{equation}\tag{HPS27}
 \mathcal H_{L}^{(r)}
   =\{(x,\nu):x\in\ell^r(\mathbb N_{\ge1}),
       \nu\in L^{\mathbb N_{\ge1}},\
       x_n\ne0\Rightarrow\nu_n=1_L\}.
\end{equation}
Addition joins labels coordinatewise and
adds the original coordinates; the scalar
$(a,\eta)\in G_L(\mathbb C)$ acts by
$(ax_n,\eta\wedge\nu_n)_n$. These laws make
it a semimodule, because each coordinate
has the HPS22 laws. The coefficient map
$(x,\nu)\mapsto(x,(\iota_L(\nu_n))_n)$
is injective by HPS25. The synchronized
map $(x,\lambda)\mapsto(x,(\lambda)_n)$
is injective and linear; its image consists
exactly of constant-label sequences.

The original shifts have the full coordinate lifts
\begin{equation}\tag{HPS28}
\begin{gathered}
 \widetilde S_m(x,\nu)=(S_mx,\nu^{(m)}),\qquad
 \nu^{(m)}_j=
    \begin{cases}\nu_{j/m},&m\mid j,\\0_L,&m\nmid j,\end{cases}\\
 \widetilde S_m^*(x,\nu)
                   =(S_m^*x,(\nu_{mn})_{n\ge1}),\\
 \widetilde U_t(u)(x,\nu)=(U_t(u)x,\nu),\qquad
 \widetilde H_t(x,\nu)=(H_tx,\nu)
                         \quad(x\in\mathcal D).
\end{gathered}
\end{equation}
They are linear for the scalar and component
laws just given. The shifts insert the actual
unsupported zero $\tau$ on the missing
coordinates; their adjoints select the stated
coordinates, with no sum over discarded labels.
The diagonal lifts preserve every label,
including at $t=0,n=1$, where $H_0e_1=0$
and a top input becomes the supported zero.
Coordinate substitution proves
$\widetilde S_m^*\widetilde S_m=I$,
$\widetilde S_m\widetilde S_k=\widetilde S_{mk}$,
and the full HPS5--8 intertwining identities.
For HPS5 the sum of the two right-hand
diagonal lifts joins identical labels.
All domain checks remain the original
amplitude checks of HPS5.
The constant-label embedding need not
intertwine the shift with its synchronized
lift: outside the multiples the former
keeps $\lambda$ and the latter inserts
$0_L$. This difference is specified next.

For any subset $A\subseteq\mathbb N_{\ge1}$
let $Q_A$ be the ordinary coordinate projection.
There are two distinct exact lifts:
\begin{equation}\tag{HPS29}
\begin{gathered}
 \mathsf Q_A(x,\nu)=(Q_Ax,(\mathbf1_A(n)\nu_n)_n),\qquad
 \mathsf Q_A^{\rm keep}(x,\nu)=(Q_Ax,\nu),\\
 \mathsf Q_A^{\rm keep}(x,\nu)
     =\mathsf Q_A(x,\nu)
                     +(0,(\mathbf1_{A^c}(n)\nu_n)_n).
\end{gathered}
\end{equation}
Here $\mathbf1_A(n)\nu_n$ means $\nu_n$
on $A$ and $0_L$ outside $A$.
Both are linear, as direct coordinate
addition and meet show. The second equals
the pointwise signed difference between
the identity lift and $\mathsf Q_{A^c}$:
on $A^c$ the amplitude cancels but its
label is $\nu_n\vee\nu_n=\nu_n$.
This proves the displayed exact zero-label
defect. The range projector
$\widetilde S_m\widetilde S_m^*$ is
$\mathsf Q_{m\mathbb N}$; the prime sieve
mask is $\mathsf Q_{\{(n,M)=1\}}$.
The product of signed complementary lifts
$\prod_{p\in\mathcal P}(I+\widehat{-1}
\mathsf Q_{p\mathbb N})$
has the same sieve amplitude and keeps
the entire original label sequence.
Indeed every factor keeps that sequence
and has amplitude $I-P_p$; finite
composition proves the assertion.
Thus its difference from the mask is
exactly the labelled-zero sequence on
the excluded prime coordinates in HPS29.

For an arbitrary infinite sequence of labels,
its trace cannot silently assume a join in $L$.
The exact universally available target is the
ordered ideal completion $\widehat L$:
its members are nonempty lower subsets
closed under finite joins. Its bottom is
$\{0_L\}$, top is $L$, meet is intersection,
and an arbitrary join is the ideal generated
by the union. The injective map
$j_L(\lambda)=\downarrow\lambda$ preserves
finite joins and meets. In particular
\begin{equation}\tag{HPS30}
\begin{gathered}
 I(\nu)=\bigvee_{n\ge1}^{\widehat L}\downarrow\nu_n
   =\{\lambda:\lambda\le\nu_{n_1}\vee\cdots\vee\nu_{n_k}
                     \text{ for some finite set of indices}\},\\
 \widehat{\operatorname{tr}}(x,\nu)
       =\left(\sum_{n\ge1}x_n,I(\nu)\right)
       \in G_{\widehat L}(\mathbb C)
                            \qquad((x,\nu)\in\mathcal H_L^{(1)}).
\end{gathered}
\end{equation}
The empty finite join is $0_L$.
The sum is absolutely convergent. If it
is nonzero, some $x_n$ is nonzero, so
$\nu_n=1_L$ and $I(\nu)=L$; thus the
displayed output satisfies the exact
carrier condition. Join-generated ideals
show $I(\nu\vee\mu)=I(\nu)\vee I(\mu)$.
Also
$I((\eta\wedge\nu_n)_n)=
\downarrow\eta\cap I(\nu)$:
if $\lambda\le\eta$ and
$\lambda\le\bigvee_{n\in F}\nu_n$, then
$\lambda\le\eta\wedge\bigvee_{n\in F}\nu_n
=\bigvee_{n\in F}(\eta\wedge\nu_n)$.
The reverse inclusion is immediate.
This proves additivity and linearity over
the scalar embedding
$(a,\eta)\mapsto(a,\downarrow\eta)$.
It also proves directly the arbitrary-join
distributivity required in $\widehat L$.

The completed support in HPS30 is principal
if and only if some finite join of the
$\nu_n$ dominates every $\nu_n$.
For if $I(\nu)=\downarrow\lambda$, its
generator $\lambda$ belongs to $I(\nu)$
and hence lies below such a finite join;
that finite join already belongs to
$\downarrow\lambda$, so it equals
$\lambda$. The converse is immediate.
Mere existence of $\sup_L\nu_n$ is not
sufficient: in $L=[0,1]$ with its usual
order and $\nu_n=1-1/n$, the generated
ideal is $[0,1)$, whereas
$\downarrow\sup_n\nu_n=[0,1]$.
Thus the actual completion and its
complete fibres have not been replaced
by a larger principal support.
If one separately chooses an $L$-valued
trace using an existing supremum, that is
a different target construction; its
principal embedding need not equal HPS30.

The coefficient tensor map commutes exactly
with this completion as well. Let
$\widehat{K_L}$ be the ordered-ideal
completion of the idempotent semiring
$K_L$, with ideal multiplication generated
by pairwise products. Define
\begin{equation}\tag{HPS31}
\begin{gathered}
 \iota^\#(I)=
    \{w:w\le\iota_L(\lambda_1)+\cdots+
                         \iota_L(\lambda_k),\ \lambda_i\in I\},\\
 \rho^\#(J)=
    \{\lambda:\lambda\le\rho_L(w_1)\vee\cdots
                            \vee\rho_L(w_k),\ w_i\in J\},\\
 \rho^\#\iota^\#=\operatorname{id}_{\widehat L},\qquad
 \iota^\# I(\nu)
           =\bigvee_n^{\widehat{K_L}}
                            \downarrow\iota_L(\nu_n).
\end{gathered}
\end{equation}
These are the ideals generated by the
respective images. Since the original
maps preserve finite sums and products,
their ideal extensions preserve arbitrary
joins and ideal products: both sides are
generated by the same images of finite
sums and pairwise products.
Applying the original retraction to
their generators proves the retraction
identity. The last formula follows by
the same generator description.
Therefore the map
$(r,I)\mapsto(r,\iota^\# I)$ is injective
and preserves all amplitude/support
operations, and the completed trace
HPS30 commutes with the complete original
coefficient embedding.

For the actual partition traces, all
nonzero spectral coefficients have top
labels. The masked full trace, each
nonempty prime stratum, and the
one-coordinate $n=1$ term all have
$I(\nu)=L$. Thus HPS10--18 pass through
HPS30 and HPS31 with the precise top
support already displayed in HPS26.
Under finite alternating exclusion the
labels on cancelled coordinates instead
remain those calculated in HPS29.
Under the growing prime sieve, its
amplitude trace tends to $(1+t)^{-s}$;
the surviving coordinate is $n=1$ and
its support remains top. The corresponding
operator formula is the actual inverse
$\mathcal M_sZ_s=g_s$, not a deletion
of $e$ or a collapse of intermediate
zero labels.

There is an exact, useful failure of commuting
coordinate-label limits with this trace.
Fix real $\beta>1$, and let $\mathcal P_j$
increase through the first $j$ primes.
Put $A_j=\{n\ge2:(n,\prod_{p\in\mathcal P_j}p)=1\}$,
and form the literal masked residual
\begin{equation}\tag{HPS32}
\begin{gathered}
 x^{(j)}_n=\mathbf1_{A_j}(n)(n+t)^{-\beta},
 \qquad
 \nu^{(j)}_n=
   \begin{cases}1_L,&n\in A_j,\\0_L,&n\notin A_j,\end{cases}\\
 a_j=\sum_{n\in A_j}(n+t)^{-\beta}>0,\qquad
 a_j\longrightarrow0,\qquad
 \widehat{\operatorname{tr}}(x^{(j)},\nu^{(j)})
                         =(a_j,L).
\end{gathered}
\end{equation}
Every finite prime set omits a prime, by
applying a prime divisor to one plus its
prime product. Thus $A_j$ is nonempty and
$a_j>0$. Each fixed integer $n\ge2$
has a prime divisor and is eventually
excluded. Dominated convergence by the
summable $(n+t)^{-\beta}$ gives
$a_j\to0$, which is also convergence
of these diagonal residuals in trace norm.
Each fixed coordinate label is eventually
$0_L$. The full coordinate limit is
therefore $(0,(0_L)_n)$, whose supported
trace is $(0,\{0_L\})$, the completed
unsupported zero. By contrast the traces
in HPS32, retaining their constant top
label while their amplitude tends to zero,
have limit $(0,L)$, the completed supported
zero. These two limits are distinct when
$L$ is nontrivial.

For the corresponding signed residual,
HPS29 keeps the original top label on
every coordinate. Its amplitude is the
same $x^{(j)}$, but its complete
coordinate limit is $(0,(1_L)_n)$ and
its supported trace is $(0,L)$.
Thus the exact discrepancy is the
all-top labelled-zero sequence discarded
by the masks. The faithful maps HPS25
and HPS31 preserve and detect it.
No topology on an arbitrary lattice
has been postulated here: the coordinate
limits of labels are eventually constant,
and each trace limit is taken in its
explicit fixed-support amplitude fibre.
This proves why trace-norm convergence of
amplitudes alone cannot justify exchanging
these two supported limiting operations.

\subsection{The actual prime operator in the critical strip}

The preceding operator product has an exact larger
domain after separating the constant coordinate.
Fix $0<r<1$ and let $\mathcal A_r$ be the disc
algebra of complex functions holomorphic on
$|t|<r$ and continuous on $|t|\le r$, with its
supremum norm and unchanged coordinate $t$.
It is complete: a uniform Cauchy sequence
converges uniformly, and Cauchy's integral
formula on each smaller circle proves that
the limit is holomorphic. Define
\begin{equation}\tag{HPS33}
\begin{gathered}
 P_0f=f(0)\mathbf1,\qquad X_r=\ker P_0,\qquad
 \mathcal A_r=\mathbb C\mathbf1\oplus X_r,\\
 (V_nf)(t)=f(t/n),\qquad
 \|V_n|_{X_r}\|=\frac1n .
\end{gathered}
\end{equation}
For $f(0)=0$, the function $f(t)/t$, with
its removable value at zero, is holomorphic.
On the boundary its modulus is at most
$\|f\|_r/r$. The maximum modulus principle,
first on smaller circles and then by
continuity at the outer boundary, gives
$|f(w)|\le(|w|/r)\|f\|_r$ for $|w|\le r$.
Apply this at $w=t/n$ to obtain the bound.
The unchanged monomial $f(t)=t$ attains
equality, proving the exact norm.
No rescaling of a programme function
has been made.

For $\operatorname{Re}s=\sigma>0$ define
operators on this actual vanishing subspace:
\begin{equation}\tag{HPS34}
\begin{gathered}
 \mathcal A_s^0=\sum_{n\ge1}n^{-s}V_n|_{X_r},\qquad
 \mathcal B_s^0=\sum_{n\ge1}\mu(n)n^{-s}V_n|_{X_r},\\
 \|\mathcal A_s^0\|,\|\mathcal B_s^0\|
                                  \le\zeta(\sigma+1),\qquad
 \mathcal B_s^0\mathcal A_s^0
          =\mathcal A_s^0\mathcal B_s^0=I_{X_r},\\
 \mathcal A_s^0
   =\lim_{\mathcal P}\prod_{p\in\mathcal P}
             (I-p^{-s}V_p|_{X_r})^{-1},\qquad
 \mathcal B_s^0
   =\lim_{\mathcal P}\prod_{p\in\mathcal P}
             (I-p^{-s}V_p|_{X_r}).
\end{gathered}
\end{equation}
HPS33 makes both absolute norm sums at most
$\sum n^{-\sigma-1}$. In the double
convolution, the norm of
$m^{-s}n^{-s}V_{mn}|_{X_r}$ is
$(mn)^{-\sigma-1}$, whose double sum is
finite. The divisor computation from
HPS15 therefore proves the inverse
identities on this larger half-plane.
Each prime factor has norm at most
$p^{-\sigma-1}<1$, so its geometric
inverse exists. Repeating the exact
finite expansions HPS16 proves the
products and limits. If $\mathcal P$
contains all primes at most $x\ge2$,
both missing-series operator errors
are at most
$\lfloor x\rfloor^{-\sigma}/\sigma$.
Logarithmic derivatives of the series
are uniformly summable on compact
subsets of $\sigma>0$. Thus both
operators and their inverses depend
holomorphically on $s$ there.

The needed scalar continuation has a
direct integral proof. For $\sigma>1$,
write $n^{-s}=s\int_n^\infty v^{-s-1}dv$
and interchange the absolutely summable
integrals to obtain
$\zeta(s)=s\int_1^\infty\lfloor v\rfloor
v^{-s-1}dv$. Consequently
\begin{equation}\tag{HPS35}
\begin{gathered}
 \zeta(s)=\frac{s}{s-1}
       -s\int_1^\infty\{v\}v^{-s-1}dv
                    \qquad(\operatorname{Re}s>0,\ s\ne1),\\
 \mathscr E_s
       =\zeta(s)P_0+\mathcal A_s^0(I-P_0)
                            \quad\hbox{on }\mathcal A_r.
\end{gathered}
\end{equation}
Here $\{v\}=v-\lfloor v\rfloor$.
The integral and its derivatives in $s$
converge uniformly on compact sets of
$\sigma>0$, since $\{v\}\le1$ and
fixed powers of $\log v$ remain
integrable against $v^{-\sigma-1}$.
This proves the meromorphic continuation,
with a simple pole of residue one at
$s=1$. The operator $\mathscr E_s$
agrees with $\mathcal E_s$ on the
disc algebra for $\sigma>1$, by
separating $f(0)$ in its convergent
series. It is meromorphic on $\sigma>0$,
and its residue at $1$ is exactly $P_0$.

The operator singularities are now completely
specified, with no assumption about where
a zero lies within this half-plane:
\begin{equation}\tag{HPS36}
\begin{gathered}
 \zeta(s)\ne0,\ s\ne1\quad\Longrightarrow\quad
 \mathscr E_s^{-1}
      =\zeta(s)^{-1}P_0+\mathcal B_s^0(I-P_0),\\
 \zeta(\rho)=0,\quad\operatorname{Re}\rho>0
 \quad\Longrightarrow\quad
 \ker\mathscr E_\rho=\mathbb C\mathbf1,\qquad
 \operatorname{ran}\mathscr E_\rho=X_r,\\
 \mathcal A_r/\operatorname{ran}\mathscr E_\rho
       \xrightarrow{\ \simeq\ }\mathbb C,\quad
                      [f]\longmapsto f(0),\\
 [P_0+\mathcal B_s^0(I-P_0)]\,\mathscr E_s
                   =\zeta(s)P_0+(I-P_0).
\end{gathered}
\end{equation}
Every formula follows by multiplying
the two direct-sum blocks, using HPS34.
In particular at a zero the $X_r$
block remains a bounded bijection,
with inverse $\mathcal B_\rho^0$.
The quotient map in the third line
is well-defined and injective because
its kernel is precisely $X_r$; it
is onto by constant functions.
The factor on the last line is a
holomorphic invertible family, its
inverse being $P_0+\mathcal A_s^0(I-P_0)$.
Thus near a zero of order $m$ the
entire singular block is exactly
the scalar vanishing factor
$(s-\rho)^m$ times a nonzero
holomorphic scalar on the constant
line. This conclusion does not
assume simple zeros.
The zero order is finite because
the scalar meromorphic function is
not identically zero, being positive
for real $s>1$.

The actual monomial and Hurwitz receivers are
\begin{equation}\tag{HPS37}
\begin{gathered}
 \mathscr E_s(t^k)=\zeta(s+k)t^k
       \quad(k\ge1,\ \operatorname{Re}s>0),\qquad
 \mathscr E_s\mathbf1=\zeta(s)\mathbf1,\\
 \mathscr E_sg_s(t)
     =\zeta(s)+\sum_{n\ge1}
                 \bigl((n+t)^{-s}-n^{-s}\bigr)
     =F_t(s),\qquad |t|\le r<1 .
\end{gathered}
\end{equation}
For the first line the $n$th operator
term is $n^{-s-k}t^k$, and its scalar
sum converges absolutely since
$\operatorname{Re}(s+k)>1$.
It is nonzero there: the absolutely
convergent scalar Möbius series in
HPS15 at $t=0$ supplies the inverse
of $\zeta(s+k)$.
For the second line, $g_s$ is
holomorphic on a neighbourhood of
the closed disc because $r<1$,
with the logarithm of $1+t$ fixed
by its value zero at $t=0$.
Its constant coordinate is $1$.
Each term of
$\mathcal A_s^0(g_s-\mathbf1)$
is exactly the displayed difference.
The bound in HPS33, or integration
of $-s(1+w)^{-s-1}$ along $w=t/n$,
makes this series uniformly
convergent on the closed disc and
on compact sets of $\sigma>0$.
For $\sigma>1$ it equals the original
trace. It consequently defines
its unique meromorphic continuation
to this disc and half-plane.
This is an explicit continuation
of the original object, not a
replacement by another function.

The differential law continues on exact
bounded domains as well. If $0<r'<r<1$,
Cauchy's formula gives
$\|f'\|_{r'}\le\|f\|_r/(r-r')$ for
$f\in\mathcal A_r$, using smaller
circles and then the limiting bound.
The differentiated series converges
in $\mathcal A_{r'}$, and yields
\begin{equation}\tag{HPS38}
\begin{gathered}
 \partial_t\mathscr E_sf
         =\mathcal E_{s+1}(\partial_tf)
                 \quad(\mathcal A_r\longrightarrow\mathcal A_{r'}),
       \qquad \operatorname{Re}s>0,\\
 \partial_tF_t(s)=-sF_t(s+1)=-DF_t(s).
\end{gathered}
\end{equation}
Here $\mathcal E_{s+1}$ is the ordinary
absolutely convergent operator on
$\mathcal A_{r'}$; its real part exceeds
one. Termwise differentiation supplies
$n^{-s-1}f'(t/n)$, whose norm sum is
at most $\zeta(\sigma+1)\|f'\|_{r'}$.
The derivative of the constant
$\zeta(s)f(0)$ vanishes, so at $s=1$
the first identity has a removable
operator-valued singularity on its
left side. In the second line,
$\partial_tg_s=-s g_{s+1}$ proves
the identity from HPS37 without
an exchange of divergent trace
series. This is the same $D$ as
HPS4, including its factor $s$.

All the preceding block maps are complex
linear and hence have the exact lifts
HPS22--25. At any zero $\rho$ in the
stated domain define the full supported
null kernel using every labelled zero
$Z_L=\{z_\lambda:\lambda\in L\}$.
The complete kernel and fibres are
\begin{equation}\tag{HPS39}
\begin{gathered}
 (\mathscr E_\rho^\uparrow)^{-1}(Z_L)
               =G_L(\mathbb C\mathbf1),\\
 (\mathscr E_\rho^\uparrow)^{-1}(\tau)
 =\begin{cases}
     \{\tau\},&0_L\ne1_L,\\
     G_L(\mathbb C\mathbf1),&0_L=1_L,
   \end{cases}\\
 (\mathscr E_\rho^\uparrow)^{-1}(v,1_L)
   =\{(\mathcal B_\rho^0v+c\mathbf1,1_L):c\in\mathbb C\}
                                   \quad(v\in X_r),\\
 (\mathscr E_\rho^\uparrow)^{-1}(0,\lambda)
                   =\{(0,\lambda)\}\quad(\lambda\ne1_L),\\
 q^\uparrow:G_L(\mathcal A_r)\longrightarrow G_L(\mathbb C),
       \qquad(f,\lambda)\longmapsto(f(0),\lambda).
\end{gathered}
\end{equation}
Outputs with amplitude outside $X_r$
have no preimage. For the displayed
top fibre use the two blocks in
HPS36. At a non-top label the
carrier condition forces the
input amplitude to be zero,
which proves its fibre assertion.
In particular a nonzero constant
is sent to $e$, which is distinct
from $\tau$ when $L$ is nontrivial.
The ordinary inverse image of the
unsupported zero does not replace
the displayed full supported
null kernel.

The map $q^\uparrow$ is a unital
semiring homomorphism, because
evaluation at zero is a unital
ring homomorphism and labels
are unchanged. Its exact fibres
are $f(0)=g(0)$ at the same label.
It has the following supported
linear quotient property.
If an $G_L(\mathbb C)$-linear map
$b:G_L(\mathcal A_r)\to M$ satisfies
$b(v,\lambda)=b(0,\lambda)$
for every admissible $v\in X_r$,
then it factors uniquely through
$q^\uparrow$. To prove this write
$f=f(0)\mathbf1+(f-f(0)\mathbf1)$
with the same admissible label
on both terms. The latter term
has the same image as $(0,\lambda)$.
This labelled zero is absorbed
by the first image: for any
semimodule element $y$ one has
$y+ey=(1+e)y=y$, and linearity
identifies this $ey$ with
$b(0,\lambda)$.
Thus $b(f,\lambda)=b(f(0)\mathbf1,\lambda)$.
The asserted factor is
$(c,\lambda)\mapsto b(c\mathbf1,\lambda)$;
constant functions prove its
uniqueness and linearity.
Conversely any such factor
has the stated property.
All these maps commute with
$\Phi$ by HPS25; its retraction
recovers their complete original
labelled fibres.

\begin{figure}[ht]
\centering
$\begin{array}{ccc}
 \mathcal H\text{ carrying }H_{t/m}
       &\xrightarrow{\quad S_m\quad}
       &\mathcal H\text{ carrying }H_t\\
 U_{t/m}(u)\ \downarrow&&\downarrow\ U_t(u)\\
 \mathcal H\text{ carrying }H_{t/m}
       &\xrightarrow{\quad m^{iu}S_m\quad}
       &\mathcal H\text{ carrying }H_t
\end{array}$
\par\medskip
$\begin{array}{ccccc}
 (g_s,\lambda)&\xrightarrow{\ \mathcal E_s^\uparrow\ }&
 (\mathcal E_sg_s,\lambda)&
 \xrightarrow{\ \mathcal M_s^\uparrow\ }&(g_s,\lambda)\\
 \Phi\downarrow&&\Phi\downarrow&&\Phi\downarrow\\
 (g_s,\iota_L\lambda)&\xrightarrow{\ (\mathcal E_s,\mathrm{id})\ }&
 (Z_s,\iota_L\lambda)&
 \xrightarrow{\ (\mathcal M_s,\mathrm{id})\ }&(g_s,\iota_L\lambda)
\end{array}$
\caption{Exact parameter-changing dilation, HPS5--8,
and the full supported operator Euler product,
HPS14--25. In the lower diagram the actual nonzero
atom requires $\lambda=1_L$; replacing its
amplitude by any general input $f$ gives the
same commuting diagram for every allowed
label, including all zero amplitudes. The
upper square keeps deformation $t$ distinct
from unitary time $u$. It does not depict a
fixed-$t$ scalar action; its exact correction
is HPS7--9. Independent coordinate masks and
all their surviving or cancelled labels are
calculated in HPS27--31.}
\end{figure}
